\documentclass[12pt,a4paper]{article}
\usepackage{cmap} % Makes the PDF copiable. See http://tex.stackexchange.com/a/64198/25761
\usepackage[T1]{fontenc}
\usepackage[brazil]{babel}
\usepackage[utf8]{inputenc}
\usepackage{amsmath}
\usepackage{amsfonts}
\usepackage{amssymb}
\usepackage{amsthm}
\usepackage{textcomp} % \degree
\usepackage{gensymb} % \degree
\usepackage[usenames,svgnames,dvipsnames]{xcolor}
\usepackage{hyperref}
\usepackage{graphicx}
\usepackage[margin=2cm]{geometry}

\hypersetup{
    colorlinks = true,
    allcolors = {blue}
}

\newcommand*\diff{\mathop{}\!\mathrm{d}}
\newcommand*\cotg{\operatorname{cotg}}

\newcommand*\prova{Prova IV}
\newcommand*\turma{TADS121-01C}
\newcommand*\disciplina{CDI0001}
\newcommand*\eu{Helder G. G. de Lima}
\newcommand*\data{04/06/2016}

\author{\eu}
\title{\prova - \disciplina}
\date{\data}

\begin{document}
\thispagestyle{empty}
\newgeometry{margin=2cm,bottom=0.5cm}
\begin{center}
\includegraphics[width=9.0cm]{marca}
\noindent\begin{tabular}{l c c r}
  \textbf{\disciplina (\turma)}
& \textbf{\prova}
& \textbf{\data}
\end{tabular}
\\ Prof. \eu\footnote{
Este é um material de acesso livre distribuído sob os termos da licença \href{https://creativecommons.org/licenses/by-sa/4.0/deed.pt_BR}{Creative Commons BY-SA 4.0}.}
\end{center}

\noindent Nome do(a) aluno(a): \rule{13cm}{0.01cm}

%\section*{Instruções}
\begin{center}\fbox{
\begin{minipage}{14cm}

{\footnotesize
\begin{itemize}
\renewcommand{\theenumi}{\Roman{enumi}}
\item Identifique-se em todas as folhas.
\item Mantenha o celular e os demais equipamentos eletrônicos desligados durante a prova.
\item Justifique cada resposta com cálculos ou argumentos baseados na teoria estudada.
\end{itemize}
}

\end{minipage}
}
\end{center}


%\section*{Questões}

\begin{enumerate}

\item Calcule apenas 5 das 7 integrais a seguir, deixando explícito como obteve todas as integrais que não forem imediatas:
\begin{enumerate}
\item (2,0) $\displaystyle \int \frac{-5x^2 +11x +4}{3 x^3-12 x^2+12 x} \diff x$
\item (2,0) $\displaystyle \int 4\arccos(2x) \diff x$
\item (2,0) $\displaystyle \int (3+\cos(x))^2\diff x$
\item (2,0) $\displaystyle \int (x^3+2x)\ln(4x) \diff x$
\item (2,0) $\displaystyle \int x(2x+6)^{11}+\cotg(x+3) \diff x$
\item (2,0) $\displaystyle \int \frac{8x+7}{\sqrt{x^2+10x+9}} \diff x$
\item (2,0) $\displaystyle \int \frac{ \cos(x) }{ e^{2x} } \diff x$
\end{enumerate}

\end{enumerate}

\begin{center}
BOA PROVA!!!
\end{center}

%\newpage
%\restoregeometry
%\section*{Respostas e observações}
%\begin{enumerate}
%\item
%\end{enumerate}
\end{document}
